\documentclass[11pt]{beamer}

\usepackage[utf8]{inputenc}
\usepackage[T1]{fontenc}
\usepackage[french]{babel}
\usepackage{xcolor}
\usepackage{graphicx}
\usepackage{listings}
\usepackage{tikz}

\usetheme{Madrid}
\definecolor{primary}{RGB}{0,102,153}
\setbeamercolor{structure}{fg=primary}
\setbeamercolor{frametitle}{fg=primary}
\setbeamercolor{alerted text}{fg=primary}

\title{Déploiement d'une Application Full-Stack}
\subtitle{Docker Compose et Kubernetes}
\author{Arif Ghazi}
\institute{ISET Tozeur - M1 DevOps}
\date{2026}

\begin{document}

\begin{frame}
  \titlepage
  %\begin{center}
  %  \includegraphics[width=2.5cm]{docker.png}\hspace{1cm}
  %  \includegraphics[width=2.5cm]{kubernetes.png}
  %\end{center}
\end{frame}

\begin{frame}{Agenda}
  \tableofcontents
\end{frame}

\section{Application et architecture}
\begin{frame}{Présentation de l'application choisie}
  \begin{itemize}
    \item Application de gestion de tâches CRUD.
    \item Interface utilisateur moderne avec React/Vite.
    \item Backend Node.js/Express exposant l'API \texttt{/api/tasks}.
    \item Stockage relationnel avec PostgreSQL.
  \end{itemize}
\end{frame}

\begin{frame}{Architecture globale}
  \begin{itemize}
    \item Frontend servi par Nginx.
    \item Backend API REST.
    \item Base de données PostgreSQL.
    \item Flux : frontend \textrightarrow backend \textrightarrow postgres.
  \end{itemize}
\end{frame}

\section{Technologies et environnement}
\begin{frame}{Technologies utilisées}
  \begin{itemize}
    \item Docker / Docker Compose
    \item Kubernetes (Kind)
    \item Node.js / Express
    \item PostgreSQL
    \item Nginx
    \item React / Vite
  \end{itemize}
\end{frame}

\begin{frame}{Environnement de travail}
  \begin{itemize}
    \item Machine de développement Windows avec Docker Desktop.
    \item Cluster Kubernetes local Kind.
    \item 1 nœud master + 2 workers.
    \item Namespace \texttt{app-dev} dédié à l'application.
  \end{itemize}
\end{frame}

\section{Phase 1 -- Docker Compose}
\begin{frame}{Architecture Docker Compose}
  \begin{itemize}
    \item Services : frontend, backend, database.
    \item Réseaux : \texttt{frontend-net}, \texttt{backend-net}.
    \item Volume persistant pour PostgreSQL.
    \item Dépendances gérées avec \texttt{depends\_on} et healthchecks.
  \end{itemize}
\end{frame}

\begin{frame}[fragile]{Best practices Docker Compose}
  \begin{itemize}
    \item Versions fixes des images.
    \item Builds multi-stage.
    \item Isolation réseau par services.
    \item Healthchecks pour l'ordre de démarrage.
    \item Persistance via volume nommé.
  \end{itemize}
\end{frame}

\begin{frame}[fragile]{Extrait \texttt{docker-compose.yml}}
\begin{lstlisting}
services:
  database:
    image: postgres:15-alpine
  backend:
    build: ./backend
  frontend:
    build: ./frontend
\end{lstlisting}
\end{frame}

\begin{frame}{Démo / validation Docker Compose}
  \begin{itemize}
    \item Commande principale : \texttt{docker compose up --build}.
    \item Vérifier avec \texttt{docker compose ps}.
    \item Tester l'interface via \texttt{http://localhost:8080}.
    \item Persistance des données après redémarrage.
  \end{itemize}
\end{frame}

\section{Phase 2 -- Kubernetes}
\begin{frame}{Architecture du cluster K3s / Kind}
  \begin{itemize}
    \item 1 master + 2 workers.
    \item Infrastructure locale pour tests.
    \item Namespace \texttt{app-dev} isolé.
    \item Images poussées via \texttt{kind load docker-image}.
  \end{itemize}
\end{frame}

\begin{frame}{Schéma global des objets Kubernetes}
  \begin{itemize}
    \item \textbf{Deployment} : frontend, backend.
    \item \textbf{Service} : NodePort pour frontend/backend.
    \item \textbf{StatefulSet} : PostgreSQL.
    \item \textbf{PVC} : stockage persistant.
    \item \textbf{ConfigMap / Secret} : configuration et mots de passe.
  \end{itemize}
\end{frame}

\begin{frame}[fragile]{Exemple ConfigMap / Secret}
\begin{lstlisting}
apiVersion: v1
kind: ConfigMap
metadata:
  name: app-config
data:
  BACKEND_URL: http://backend-service:4000
---
apiVersion: v1
kind: Secret
metadata:
  name: db-secret
type: Opaque
data:
  POSTGRES_PASSWORD: cGFzc3dvcmQ=
\end{lstlisting}
\end{frame}

\begin{frame}{Persistent Volume & PVC}
  \begin{itemize}
    \item Volume persistant pour PostgreSQL.
    \item PVC alloue du stockage stable.
    \item Permet de conserver les données après redémarrage.
  \end{itemize}
\end{frame}

\begin{frame}[fragile]{Extrait Deployment / Service YAML}
\begin{lstlisting}
kind: Deployment
metadata:
  name: backend
spec:
  replicas: 2
  template:
    metadata:
      labels:
        app: backend
---
kind: Service
metadata:
  name: backend-service
spec:
  type: NodePort
  ports:
    - port: 4000
\end{lstlisting}
\end{frame}

\begin{frame}{Best practices Kubernetes appliquées}
  \begin{itemize}
    \item Readiness / liveness probes.
    \item Labels cohérents et selectors clairs.
    \item Secrets pour les informations sensibles.
    \item StatefulSet pour la base de données.
    \item NetworkPolicy pour l'isolation réseau.
  \end{itemize}
\end{frame}

\section{Tests et validation}
\begin{frame}{Tests et validation}
  \begin{itemize}
    \item Vérifier les pods avec \texttt{kubectl get pods -n app-dev}.
    \item Contrôler les services avec \texttt{kubectl get svc -n app-dev}.
    \item Tester le frontend avec port-forward.
    \item Valider les endpoints API CRUD.
  \end{itemize}
\end{frame}

\begin{frame}[fragile]{Test réseau / NetworkPolicy}
\begin{lstlisting}
kubectl run test-frontend -n app-dev --restart=Never \
  --rm --image=curlimages/curl \
  -- sh -c "curl --max-time 5 -sS http://postgres-service:5432 >/dev/null 2>&1; echo EXIT_CODE=$?"
\end{lstlisting}
\end{frame}

\begin{frame}{Résilience et disponibilité}
  \begin{itemize}
    \item Suppression manuelle d'un pod frontend/backend.
    \item Kubernetes recrée automatiquement le pod.
    \item Pod redevenu \texttt{Ready} en environ 30 secondes.
  \end{itemize}
\end{frame}

\section{Difficultés et solutions}
\begin{frame}{Difficultés rencontrées}
  \begin{itemize}
    \item Synchronisation database/backend.
    \item Chargement des images dans Kind.
    \item Configuration initiale des NetworkPolicy.
    \item Port-forward occupé par processus actif.
  \end{itemize}
\end{frame}

\begin{frame}{Solutions apportées}
  \begin{itemize}
    \item Healthchecks et \texttt{depends\_on}.
    \item Chargement explicite des images avec Kind.
    \item Ajout de policies Ingress/Egress.
    \item Arrêt des processus conflictuels avant test.
  \end{itemize}
\end{frame}

\section{Conclusion}
\begin{frame}{Conclusion et améliorations futures}
  \begin{itemize}
    \item Projet opérationnel en Docker Compose et Kubernetes.
    \item Isolation réseau et résilience validées.
    \item Améliorations : CI/CD, monitoring, GitOps, production.
  \end{itemize}
\end{frame}

\begin{frame}{Merci}
  \begin{center}
    \Large Questions ?
  \end{center}
\end{frame}

\end{document}
